%=======================================================================
%  CHAPTER 7 - DISCRETE FOURIER TRANSFORM AND FFT  (7 hrs)
%  Every DFT value, butterfly stage, circular convolution and complexity
%  figure quoted here is produced by dft.py, which self-tests DIT and DIF
%  against the direct DFT on nine sequences before it prints, and is
%  asserted again by verify.py ch7.
%  Sources: SPP Sir deck 7 "DFT", 45 slides, and deck 8 "FFT", 26 slides,
%           both EXPORTED THIS SESSION with tools/ppt_export.py. These are
%           the syllabus-aligned source and the figure supply:
%             deck 7 p7-14  DFS -> DFT, the analysis and synthesis pair
%             deck 7 p15-31 the properties, ending in Oppenheim Table 8.2
%             deck 7 p32-45 linear convolution via the DFT, the aliasing
%                           analysis, and zero padding
%             deck 8 p5-16  decimation in time, ending in the flow graph
%             deck 8 p17-26 decimation in frequency, same
%           Chapterwise Ch4 "The DFT and FFT" and Ch7 "Discrete Fourier
%           transform" are two different authors on the same material.
%           Notes by BB Sir Chapter 7.1b (49 pp, DFT) and 7.2a / 7.2b
%           (FFT). NOTE the handoff claimed these seven files are near
%           duplicates; they share a title slide and nothing else, one
%           common page signature across six files.
%  Numericals live in ch7-num.tex.
%=======================================================================
\section{Discrete Fourier Transform and FFT}

{\color{sub}\footnotesize 7 hrs \gap$\cdot$\gap \mk{99 questions across the 45 papers}
\gap$\cdot$\gap worth about \mk{14 of 80 marks} \gap$\cdot$\gap 3 syllabus sub-topics
\gap$\cdot$\gap \mk{the heaviest chapter in the subject}.}

\lead{Read this first:}
\begin{itemize}
  \item \mk{Three questions, and only three}, in 44 of the 45 papers:
  \begin{itemize}
    \item \textbf{Compute an $N$-point DFT by DIT or DIF} \tS{43}, 340 marks. Almost
          always \mk{8-point}, almost always "draw the butterfly".
    \item \textbf{Circular convolution} \tS{39}, 253 marks, in three disguises that are
          all the same calculation (\S7.5).
    \item \textbf{Bookwork on the DFT} \tS{17}, 52 marks: why DFT, why FFT, the
          properties, zero padding.
  \end{itemize}
  \item \mk{The FFT is not a new transform.} It is the DFT, computed by a rearrangement
        that costs $\tfrac{N}{2}\log_2N$ multiplications instead of $N^2$. Anyone who
        says "the FFT of $x[n]$" means its DFT.
  \item \mk{DIT and DIF give identical answers.} They differ only in where the
        bit-reversal sits and where the twiddle factor is applied. If a paper does not
        say which, either is a full-marks answer.
  \item \textbf{Everything rests on two identities}, and nothing else:
  \item[] \[ W_N=e^{-j2\pi/N},\qquad
             W_N^{\,k+N/2}=-W_N^{\,k},\qquad W_N^{\,k+N}=W_N^{\,k} \]
        The first halves the work at every stage; the second is what makes the
        recursion close.
  \item \mk{A sequence shorter than $N$ is zero padded.} Eleven papers print six or
        seven values for an 8-point transform. That is not a misprint to fix, it is
        the zero-padding they ask you to define elsewhere in the same paper.
\end{itemize}

\hr

\T{7.1 \quad What the DFT is, and why it exists}

\Q{Why do we need the DFT / FFT although we have the DTFT? Differentiate between the
DFT and the DTFT}{\m{1+5}\m{2+4}\m{6}\m{2}}
{\color{sub}\footnotesize \tS{13} \yr{\textbf{79 Bh}, \textbf{78 Bh}, \textbf{71 Ch},
\textbf{81 Bh}, \textbf{\texttt{79 Ch}}, \textbf{\texttt{70 Bh}}, 79 Ba, 80 Ba,
\texttt{74 Ma}} \gap$\cdot$\gap \mk{asked in nine papers, always the same two reasons}}

\sbs{%
\lead{Why the DTFT is not enough}
\begin{itemize}
  \item The DTFT $X(e^{j\omega})=\sum_n x[n]e^{-j\omega n}$ is a \textbf{continuous}
        function of $\omega$.
  \item \mk{A computer cannot store a continuous function}, and it cannot sum an
        infinite series. Neither end of the DTFT is computable.
  \item The DFT fixes both at once: \textbf{finite} sequence in, \textbf{finite} set of
        samples out.
  \item[] \[ X[k]=\sum_{n=0}^{N-1}x[n]W_N^{\,nk},\qquad
             x[n]=\frac1N\sum_{k=0}^{N-1}X[k]W_N^{-nk} \]
  \item \mk{The DFT is the DTFT sampled at $N$ equally spaced points},
        $\omega_k=2\pi k/N$. Nothing more.
  \item The price: the DFT implicitly treats $x[n]$ as \textbf{one period of a periodic
        sequence}. Every oddity in this chapter, circular shift, circular convolution,
        wrap-around, follows from that one fact.
\end{itemize}}{%
\lead{DFT against DTFT, as a table}
\par\smallskip
\begin{tabular}{@{}L{2.2cm}L{2.8cm}L{2.8cm}@{}}
\toprule
\hrow \thead{} & \thead{DTFT} & \thead{DFT} \\
Input & any length, even infinite & \textbf{finite}, $N$ points \\
Output & continuous in $\omega$ & \textbf{discrete}, $N$ points \\
Periodic in & $\omega$, period $2\pi$ & $k$, period $N$ \\
Computable & \mk{no} & \mk{yes} \\
Assumes & nothing & $x[n]$ is periodic with period $N$ \\
Convolution & linear & \textbf{circular} \\
\bottomrule
\end{tabular}
\lead{And why the FFT}
\begin{itemize}
  \item The DFT as defined costs $N^2$ complex multiplications. \mk{The FFT is an
        algorithm for the same numbers} at $\tfrac{N}{2}\log_2N$.
  \item For $N=8$: $64$ against $12$. For $N=1024$: $1\,048\,576$ against $5120$, a
        \mk{speed-up of 204.8}. Quote a number; it is worth a mark.
\end{itemize}}

\Q{How does the FFT reduce computational complexity compared with direct computation of
the DFT?}{\m{5}\m{2}\m{4+2}}
{\color{sub}\footnotesize \tS{12} \yr{the same nine papers, plus
\textbf{\texttt{74 Bh}} \m{2}, \texttt{73 Ma} \m{2}, \textbf{\texttt{79 Ch}} \m{4+2}}}

\sbs{%
\lead{The argument, in four lines}
\begin{enumerate}
  \item Direct DFT: $N$ outputs, each a sum of $N$ products.
        \mk{$N^2$ complex multiplications and $N(N-1)$ additions.}
  \item Split $x[n]$ into two half-length sequences. Two $\tfrac{N}{2}$-point DFTs cost
        $2(\tfrac{N}{2})^2=\tfrac{N^2}{2}$, \textbf{half the work}, plus $\tfrac{N}{2}$
        twiddle multiplications to recombine.
  \item \mk{Repeat the split.} $N=2^m$ allows it $m=\log_2N$ times, until only
        2-point DFTs remain, and a 2-point DFT needs \textbf{no multiplication at all}
        ($W_2^0=1$).
  \item Each of the $\log_2N$ stages has $\tfrac{N}{2}$ butterflies and each butterfly
        costs one multiplication:
\end{enumerate}
\[ \boxed{\ \tfrac{N}{2}\log_2N \ \text{multiplications},\quad
   N\log_2N \ \text{additions}\ } \]}{%
\lead{The numbers to quote}
\par\smallskip
\begin{tabular}{@{}L{1.5cm}L{2.2cm}L{2.0cm}L{1.6cm}@{}}
\toprule
\hrow \thead{$N$} & \thead{Direct $N^2$} & \thead{FFT $\frac N2\log_2N$}
& \thead{Speed-up} \\
$8$ & $64$ & $12$ & $5.33$ \\
$32$ & $1024$ & $80$ & $12.8$ \\
$256$ & $65\,536$ & $1024$ & $64$ \\
$1024$ & $1\,048\,576$ & $5120$ & \mk{$204.8$} \\
\bottomrule
\end{tabular}
\begin{itemize}
  \item \mk{The saving grows with $N$}, which is the point: it is $2N/\log_2N$, not a
        constant factor.
  \item \textbf{Radix-2 needs $N$ a power of 2.} If it is not, zero pad up to one. That
        is the second reason zero padding exists.
  \item The two properties being exploited are \textbf{periodicity}
        $W_N^{k+N}=W_N^k$ and \textbf{symmetry} $W_N^{k+N/2}=-W_N^k$. Name them: they
        are usually worth a mark each.
\end{itemize}}

\hr

\T{7.2 \quad Properties of the DFT}

\Q{State and prove the circular convolution property of the DFT; write about the
multiplication and convolution properties}{\m{2+2+4}\m{6}\m{1}}
{\color{sub}\footnotesize \tF{4} \yr{\textbf{70 Ch} \m{2+2+4}, \textbf{69 Bh} \m{6},
75 Ash \m{1}, \textbf{\texttt{73 Bh}} \m{1}} \gap$\cdot$\gap \mk{the proof is five lines}}

\sbsr{0.52}{%
\lead{The property, and its proof}
\begin{itemize}
  \item \textbf{Statement.} If $x_3[n]=x_1[n]\circledast x_2[n]$, the \emph{circular}
        convolution, then
  \item[] \[ X_3[k]=X_1[k]\,X_2[k] \]
        \mk{Circular convolution in time is multiplication in frequency.}
  \item \textbf{Proof.} Write the circular convolution and transform it:
  \item[] \[ X_3[k]=\sum_{n=0}^{N-1}\Big(\sum_{m=0}^{N-1}x_1[m]\,
             x_2[((n-m))_N]\Big)W_N^{\,nk} \]
  \item Swap the sums and put $\ell=n-m$, so $n=\ell+m$ and $W_N^{nk}=W_N^{mk}W_N^{\ell k}$:
  \item[] \[ =\sum_{m=0}^{N-1}x_1[m]W_N^{\,mk}
             \sum_{\ell=0}^{N-1}x_2[((\ell))_N]W_N^{\,\ell k}
             =X_1[k]X_2[k] \]
  \item \mk{The modulo-$N$ index is what makes the inner sum come out as $X_2[k]$
        exactly}, with no leftover terms. That is the whole proof, and the step worth
        pointing at.
  \item \textbf{The dual (multiplication property)}: multiplying in time is
        \emph{circular convolution divided by $N$} in frequency,
        $x_1[n]x_2[n]\longleftrightarrow\frac1N X_1[k]\circledast X_2[k]$.
        \mk{Do not forget the $1/N$.}
\end{itemize}}{%
\lead{The full property table \mk{(Oppenheim Table 8.2)}}
\par\smallskip
{\footnotesize
\begin{tabular}{@{}L{3.3cm}L{3.5cm}@{}}
\toprule
\hrow \thead{Sequence, length $N$} & \thead{$N$-point DFT} \\
$x[n]$ & $X[k]$ \\
$ax_1[n]+bx_2[n]$ & $aX_1[k]+bX_2[k]$ \\
$X[n]$ & $Nx[((-k))_N]$ \\
$x[((n-m))_N]$ & $W_N^{\,km}X[k]$ \\
$W_N^{-\ell n}x[n]$ & $X[((k-\ell))_N]$ \\
$\sum_m x_1[m]x_2[((n-m))_N]$ & $X_1[k]X_2[k]$ \\
$x_1[n]x_2[n]$ & $\frac1N\sum_\ell X_1[\ell]X_2[((k-\ell))_N]$ \\
$x^*[n]$ & $X^*[((-k))_N]$ \\
$x[n]$ real & $X[k]=X^*[((-k))_N]$ \\
\bottomrule
\end{tabular}}
\begin{itemize}
  \item \mk{Row 4 is the circular shift}: a shift in time is a \emph{phase} in
        frequency and nothing else. The magnitude spectrum does not move.
  \item \mk{Row 3 is duality}, and it is the whole of the 67 Mng question in \S7.3.
  \item The last row is the free check on every answer in this chapter:
        \textbf{a real sequence has a conjugate-symmetric spectrum}, so
        $X[N-k]=X^*[k]$. Compute half the values and mirror the rest.
\end{itemize}}

\Q{Define zero padding, and Parseval's theorem}{\m{1}\m{2}}
{\color{sub}\footnotesize \tS{7} \yr{\textbf{74 Ch}, \textbf{74 Ash}, \textbf{70 Ch},
\textbf{\texttt{81 Ch}}, \textbf{\texttt{75 Bh}}, \textbf{\texttt{74 Bh}} \m{1};
\textbf{\texttt{77 Ch}} \m{2}} \gap$\cdot$\gap \mk{eight free marks}}

\sbs{%
\lead{Zero padding}
\begin{itemize}
  \item \textbf{Appending zeros} to a sequence before taking its DFT, to reach a
        chosen length $N$.
  \item \textbf{Three reasons}, and papers accept any of them:
  \begin{itemize}
    \item to reach a \mk{power of 2} so a radix-2 FFT can run;
    \item to make $N\ge L+M-1$ so a \mk{circular convolution equals the linear one}
          (\S7.5, and this is the reason the question sits beside those problems);
    \item to \mk{interpolate the spectrum}: more samples of the same underlying DTFT.
  \end{itemize}
  \item \mk{Zero padding does NOT add information and does NOT improve resolution.}
        It samples the same DTFT more finely. Saying otherwise loses the mark.
\end{itemize}}{%
\lead{Parseval's theorem}
\begin{itemize}
  \item Energy is the same computed either side of the transform:
  \item[] \[ \sum_{n=0}^{N-1}|x[n]|^2=\frac1N\sum_{k=0}^{N-1}|X[k]|^2 \]
  \item \mk{Note the $1/N$}, which is on the frequency side, matching the IDFT.
  \item Useful as an arithmetic check: compute $\sum|x[n]|^2$ from the given sequence in
        two seconds, then confirm the $X[k]$ you just worked out.
\end{itemize}
\lead{Linear against circular convolution \mk{(75 Ch asks for it directly)}}
\begin{itemize}
  \item \textbf{Linear}: indices run freely, result has length $L+M-1$, and it is what a
        real LTI system does.
  \item \textbf{Circular}: indices are taken $\bmod\ N$, result has length $N$, and it
        is what multiplying two DFTs gives you.
  \item \mk{They agree exactly when $N\ge L+M-1$}, and otherwise the circular one is the
        linear one \textbf{wrapped round and added to itself}.
\end{itemize}}

\hr

\T{7.3 \quad The IDFT from a forward DFT}

\Q{Using the duality property, how can a forward DFT routine be used to compute the
IDFT?}{\m{8}\m{2+6}}
{\color{sub}\footnotesize \tP{2} \yr{\textbf{67 Mng} \m{8},
\textbf{\texttt{72 Ash}} \m{2+6}} \gap$\cdot$\gap \mk{8 marks for four lines of algebra}}

\sbs{%
\lead{The conjugate trick, which is the expected answer}
\begin{enumerate}
  \item Start from the IDFT and conjugate both sides:
  \item[] \[ x[n]=\frac1N\sum_k X[k]W_N^{-nk}
             \ \Longrightarrow\
             x^*[n]=\frac1N\sum_k X^*[k]W_N^{\,nk} \]
  \item \mk{The right-hand sum is a forward DFT of $X^*[k]$.} So:
  \item[] \[ \boxed{\,x[n]=\frac1N\Big(\mathrm{DFT}\big\{X^*[k]\big\}\Big)^{*}\,} \]
  \item \textbf{Three steps}: conjugate the input, run the same FFT, conjugate the
        output and divide by $N$.
  \item \mk{No second routine is needed}, which is the point of the question: the same
        butterfly hardware computes both directions.
\end{enumerate}}{%
\lead{The duality route, if the paper insists}
\begin{itemize}
  \item Duality says $X[n]\xrightarrow{\ \mathrm{DFT}\ }Nx[((-k))_N]$.
  \item So feeding the \emph{spectrum} $X[k]$ into a forward DFT returns the original
        sequence, scaled by $N$ and \textbf{index reversed modulo $N$}:
  \item[] \[ x[n]=\frac1N\,\mathrm{DFT}\big\{X[k]\big\}\Big|_{\,k\to-n \bmod N} \]
  \item Index reversal modulo $N$ leaves entry $0$ alone and reverses the rest, so
        $\{a,b,c,d\}\to\{a,d,c,b\}$. \mk{That single sentence is what the 8 marks
        are for.}
  \item Both routes are correct. The conjugate one is easier to defend and is what
        72 Ash's worked answer needs (\S2 of the numericals).
\end{itemize}}

\hr

\T{7.4 \quad The FFT: decimation in time and in frequency}

\Q{How do you use the DIT-FFT algorithm to compute the 8-point DFT of a sequence?
Explain with the flow graph}{\m{6}\m{4}}
{\color{sub}\footnotesize \tS{43} \yr{\texttt{70 Ma} \m{6}, \textbf{69 Bh} \m{4}
(4-point), and the 41 computation papers in \S7.6}}

\sbsr{0.46}{%
\lead{Decimation in time, derived}
\begin{itemize}
  \item Split the sum by \textbf{even and odd $n$}:
  \item[] \[ X[k]=\sum_{r}x[2r]W_N^{\,2rk}+\sum_{r}x[2r+1]W_N^{\,(2r+1)k} \]
  \item $W_N^{2}=W_{N/2}$, so both sums are $\tfrac N2$-point DFTs:
  \item[] \[ \boxed{\,X[k]=G[k]+W_N^{\,k}H[k]\,} \]
        with $G$ from the even samples and $H$ from the odd ones.
  \item The other half comes free from the symmetry:
  \item[] \[ X[k+\tfrac N2]=G[k]-W_N^{\,k}H[k] \]
  \item \mk{One multiplication serves two outputs.} That pair is the butterfly, and it
        is the entire saving.
  \item Recurse $\log_2N$ times. Splitting by even and odd repeatedly leaves the input
        in \textbf{bit-reversed order}; the output comes out natural.
\end{itemize}}{%
\lead{The butterfly}
\figT{c7_butterfly.png}
{\color{sub}\footnotesize Left, the butterfly as first drawn: two twiddles. Right, the
same butterfly after $W_N^{\,r+N/2}=W_N^{\,r}W_N^{\,N/2}=-W_N^{\,r}$ is used, so one
multiplication and a sign. \mk{Always draw the right-hand form.}}
\begin{itemize}
  \item Every stage is $\tfrac N2$ of these; there are $\log_2N$ stages.
  \item \mk{The $-1$ on the lower branch is not decoration.} Omitting it is the single
        most common way to lose the marks on a flow graph.
\end{itemize}}

\lead{The 8-point DIT flow graph \mk{(draw this one)}}
\figC{c7_dit_flow.png}{0.78\textwidth}
{\color{sub}\footnotesize \mk{Input bit-reversed} $\{x[0],x[4],x[2],x[6],x[1],x[5],x[3],x[7]\}$,
output natural. Three stages. Twiddles $W_N^0$ in stage 1; $W_N^0,W_N^2$ in stage 2;
$W_N^0,W_N^1,W_N^2,W_N^3$ in stage 3. \textbf{The twiddle count doubles each stage and
the butterfly span doubles with it.}}

\Q{Draw the flow graph for the 8-point decimation in frequency algorithm}{\m{8}\m{2+6}}
{\color{sub}\footnotesize \tS{20} \yr{the DIF half of \S7.6: \textbf{80 Bh},
\textbf{76 Ch}, 76 Ash, 79 Ba, \textbf{73 Ch}, \textbf{69 Ch}, 74 Ash,
\textbf{\texttt{70 Bh}}, \textbf{\texttt{81 Ch}}, \textbf{\texttt{78 Ch}},
\textbf{\texttt{77 Ch}}, \textbf{\texttt{75 Bh}}, \texttt{72 Ma}, \textbf{68 Bh},
82 Ba, 80 Ba, \textbf{70 Ch}, \textbf{74 Ch}}}

\sbsr{0.46}{%
\lead{Decimation in frequency, derived}
\begin{itemize}
  \item This time split the \textbf{input} into first and second halves, and the
        \textbf{output} into even and odd $k$:
  \item[] \[ g[n]=x[n]+x[n+\tfrac N2],\qquad
             h[n]=\big(x[n]-x[n+\tfrac N2]\big)W_N^{\,n} \]
  \item[] \[ X[2r]=\mathrm{DFT}_{N/2}\{g[n]\},\qquad
             X[2r+1]=\mathrm{DFT}_{N/2}\{h[n]\} \]
  \item \mk{The twiddle multiplies the DIFFERENCE, after the subtraction.} In DIT it
        multiplies before the add. That is the only real difference between the two.
  \item Recurse $\log_2N$ times. Input stays natural; \textbf{output comes out
        bit-reversed} and has to be unscrambled.
\end{itemize}
\lead{DIT against DIF, the whole comparison}
\par\smallskip
{\footnotesize
\begin{tabular}{@{}L{1.9cm}L{2.0cm}L{2.0cm}@{}}
\toprule
\hrow \thead{} & \thead{DIT} & \thead{DIF} \\
Splits & input, even/odd $n$ & output, even/odd $k$ \\
Input order & \mk{bit-reversed} & natural \\
Output order & natural & \mk{bit-reversed} \\
Twiddle & \textbf{before} the add & \textbf{after} the subtract \\
Stages & $\log_2N$ & $\log_2N$ \\
Multiplies & $\frac N2\log_2N$ & $\frac N2\log_2N$ \\
Answer & \multicolumn{2}{l}{\mk{identical}} \\
\bottomrule
\end{tabular}}}{%
\lead{Bit reversal, which is all either algorithm scrambles}
\begin{itemize}
  \item \mk{The same permutation serves both}: $0,4,2,6,1,5,3,7$ for $N=8$. It is its
        own inverse, so there is only one order to memorise.
  \item Get it by writing the index in binary and reading it \textbf{backwards}:
        $3=011\to110=6$, $1=001\to100=4$.
  \item \mk{Exactly one end is scrambled}, never both and never neither. DIT scrambles
        the input, DIF the output.
  \item For $N=4$ it is $0,2,1,3$; for $N=16$, $0,8,4,12,2,10,6,14,1,9,5,13,3,11,7,15$.
  \item The scrambling is not overhead you can avoid. It is what lets every butterfly
        write its result back over its own inputs, so the whole transform runs in
        \textbf{one array with no scratch space}.
\end{itemize}}

% the figure is about 20 baselines tall and its caption four more; without a
% reserve the caption orphans onto the next page, away from the graph it names
\Needspace{23\baselineskip}
\lead{The 8-point DIF flow graph \mk{(the other one to be able to draw)}}
\figC{c7_dif_flow.png}{0.66\textwidth}
{\color{sub}\footnotesize Input natural, \mk{output bit-reversed}:
$X[0],X[4],X[2],X[6],X[1],X[5],X[3],X[7]$ read down the right edge. The dashed boxes
show the recursion, the $N/2=4$-point DFTs inside the whole transform and the
$N/4=2$-point DFTs inside those. \textbf{The twiddles run $W_N^0$ to $W_N^3$ in stage 1}
and halve their range at each stage, the mirror image of DIT, where they start at
$W_N^0$ alone and double.}

\hr

% \T reserves 12 baselines, which covers the band and its question heading but
% not the year-tag line, the note and the unbreakable \sbs that follow. Move
% the whole group together.
\Needspace{24\baselineskip}
\T{7.5 \quad Circular convolution, and the three ways it is asked}

\Q{Compute the circular convolution of two sequences}{\m{5}\m{6}\m{7}\m{2+6}}
{\color{sub}\footnotesize \tS{39} \yr{15 papers ask it directly, 12 ask for
$x_3=\mathrm{IDFT}(X_1X_2)$, 4 ask for a linear convolution done circularly. All 31
are one calculation.}}

{\color{sub}\footnotesize In one line: \mk{circular convolution is ordinary convolution
with the index taken modulo $N$}, so the tail that would have run off the end wraps
round and adds onto the front instead. Everything below is that sentence applied.}

\sbs{%
\lead{The three disguises}
\begin{enumerate}
  \item \textbf{"Compute the circular convolution of $x_1$ and $x_2$."} Direct.
  \item \textbf{"Find $x_3[n]$ if $X_3[k]=X_1[k]X_2[k]$."} By the property in \S7.2 this
        \mk{IS the circular convolution}, so do the same arithmetic and do not compute a
        single DFT.
  \item \textbf{"Find the linear convolution using circular convolution."} Zero pad both
        to $N\ge L+M-1$, then the same arithmetic again.
\end{enumerate}
{\color{sub}\footnotesize \mk{Recognising that (2) needs no transform is worth several
minutes in the exam.} The long route, three DFTs and an IDFT, earns the same marks and
takes four times as long. Say in one line that you are using the property, then convolve.}}{%
\lead{The matrix method, which is the fastest by hand}
\begin{itemize}
  \item Write the \textbf{circulant} of one sequence: first column is $x_2$, and each
        next column is the previous rotated \emph{down} by one.
  \item[] \[ y[n]=\sum_{m=0}^{N-1}x_1[m]\,x_2[((n-m))_N] \]
  \item Multiply by $x_1$ as a column vector. \mk{Each row of the product is one output
        sample}, so the whole answer falls out of one $N\times N$ table.
  \item Check: \mk{$\sum y[n]=\big(\sum x_1\big)\big(\sum x_2\big)$}, always. It is
        $Y[0]=X_1[0]X_2[0]$ in disguise, it costs two additions, and it catches almost
        every slip.
\end{itemize}}

\sbs{%
\lead{Choosing $N$}
\begin{itemize}
  \item If the paper names $N$, use it, and pad both sequences up to it.
  \item If not, $N=\max(L,M)$ for a circular answer, and $N=L+M-1$ for a linear one.
  \item \mk{$N$ smaller than $L+M-1$ wraps}, and the wrap is not an error: it is the
        correct circular answer to the question that was asked.
  \item \mk{$N$ larger than $L+M-1$ only appends zeros} to the result. Harmless, and
        sometimes needed to reach a power of 2.
\end{itemize}}{%
\lead{The concentric-circle picture}
\begin{itemize}
  \item The same arithmetic, drawn: $x_1$ clockwise on an inner circle, $x_2$
        anticlockwise on an outer one, multiply the aligned pairs and sum, then rotate
        one circle a step for the next output.
  \item \mk{Use the matrix instead.} The circles are slower, easier to misalign, and no
        paper in the 45 demands them.
\end{itemize}}

\hr

\T{7.6 \quad Last-minute recall}

\sbs{%
\lead{Say these without thinking}
\begin{itemize}
  \item $X[k]=\sum_{n=0}^{N-1}x[n]W_N^{nk}$, $W_N=e^{-j2\pi/N}$, IDFT carries $1/N$
        and $W_N^{-nk}$.
  \item $W_4$: $1,-j,-1,+j$. $W_8^1=\tfrac{1-j}{\sqrt2}$, $W_8^2=-j$,
        $W_8^3=\tfrac{-1-j}{\sqrt2}$.
  \item $W_N^{k+N/2}=-W_N^k$ and $W_N^{k+N}=W_N^k$.
  \item Bit reversal for $N=8$: \mk{$0,4,2,6,1,5,3,7$}.
  \item FFT cost $\tfrac N2\log_2N$ multiplies, $N\log_2N$ adds; $N=1024$ gives
        $204.8\times$.
  \item DIT: bit-reversed in, natural out, twiddle before the add.
        DIF: natural in, bit-reversed out, twiddle after the subtract.
  \item $X_3[k]=X_1[k]X_2[k]\Leftrightarrow$ \textbf{circular} convolution.
  \item Linear $=$ circular when $N\ge L+M-1$.
  \item $X[0]=\sum x[n]$, and for real $x$, $X[N-k]=X^*[k]$.
  \item Parseval: $\sum|x[n]|^2=\frac1N\sum|X[k]|^2$.
\end{itemize}}{%
\lead{Mistakes that cost marks}
\begin{itemize}
  \item \mk{Forgetting the $-1$} on the lower butterfly branch.
  \item Feeding DIT in natural order, or reading DIF's output in natural order.
        \textbf{Exactly one of the two ends is scrambled, never both and never neither.}
  \item Using $e^{+j2\pi/N}$ in the forward transform. \mk{The forward sign is minus.}
  \item Computing three DFTs and an IDFT for a question that only wanted a circular
        convolution.
  \item Treating a 6 or 7 entry sequence as an error. \mk{Pad it and say you padded it.}
  \item Claiming zero padding improves frequency resolution. It does not.
  \item Dropping the $1/N$ in the IDFT, or in the multiplication property.
  \item Giving a circular convolution of length $L+M-1$. \mk{It has length $N$.}
  \item Forgetting that $X[k]$ repeats: $X[5]$ of a 4-point DFT is $X[1]$, which is
        exactly what 75 Ch is testing.
\end{itemize}}
